% ============================================================================
% Paper 12: How Well Do PEFT Methods Adapt Geospatial Foundation Models
% to Heterogeneous Inputs? A Systematic Evaluation on Prithvi-100M
% ============================================================================
% Target: GIScience & Remote Sensing / CVPR EarthVision Workshop
% ============================================================================
\documentclass{article}

\usepackage{amsmath,amssymb}
\usepackage{booktabs}
\usepackage{graphicx}
\usepackage{multirow}
\usepackage{caption}
\usepackage{subcaption}
\usepackage{hyperref}
\usepackage{color}
\usepackage{natbib}
\usepackage{tikz}
\usetikzlibrary{positioning,arrows.meta,shapes.geometric,fit,calc}
\bibliographystyle{plainnat}

\begin{document}

\makeatletter
\@ifundefined{ANONYMIZE}{
  \title{How Well Do PEFT Methods Adapt Geospatial Foundation Models\\to Heterogeneous Inputs?\\A Systematic Evaluation on Prithvi-100M}
  \author{Ning Zhou$^{a}$ and Xiang Jing$^{a}$$^{\ast}$\thanks{$^\ast$Corresponding author. Email: jingxiang@pku.edu.cn}\\
  $^{a}${\em School of Software and Microelectronics, Peking University, Beijing 100871, China}}
}{
  \title{How Well Do PEFT Methods Adapt Geospatial Foundation Models to Heterogeneous Inputs? A Systematic Evaluation on Prithvi-100M}
  \author{Anonymous authors}
}
\makeatother

\date{}
\maketitle

% ============================================================================
\begin{abstract}
% ============================================================================

Geospatial foundation models (GeoFMs) such as Prithvi-100M are pre-trained on fixed spectral band configurations, yet real-world deployment requires adaptation to heterogeneous inputs---commercial high-resolution optical imagery with 4 bands, SAR with 2 bands, or arbitrary sensor combinations. Current practice handles this mismatch through zero-padding, but no systematic study has evaluated how different parameter-efficient fine-tuning (PEFT) methods perform under such conditions. We present the first comprehensive benchmark of five PEFT strategies---Linear Probing, BitFit, LoRA, Houlsby Adapter, and a novel input-stage residual adapter (GeoAdapter)---applied to a frozen Prithvi-100M backbone across five modality configurations on EuroSAT (75 experiments, 3 seeds each). Our results reveal three findings with practical implications: (1)~LoRA, widely considered the most effective PEFT method, completely fails on Prithvi's fused-QKV attention architecture ($\Delta$OA~$<$~0.001 vs.\ Linear Probe); (2)~Houlsby adapters dominate all configurations with +16--23\% OA improvement and can implicitly utilize zero-padded channels for cross-modal transfer (rgb\_sar outperforms rgb by +1.1\%); (3)~input-stage modality adaptation with a lightweight residual adapter consistently underperforms zero-padding, demonstrating that the modality gap is better addressed through backbone-internal adaptation than input-stage projection. We release our evaluation platform, Geo-MLOps, as open-source software to facilitate reproducible PEFT research on geospatial foundation models.

\end{abstract}

\begin{keywords}
parameter-efficient fine-tuning; geospatial foundation model; Prithvi; heterogeneous inputs; benchmark; negative results
\end{keywords}

% ============================================================================
\section{Introduction}
\label{sec:intro}
% ============================================================================

Geospatial foundation models (GeoFMs) have rapidly advanced the state of the art in Earth observation tasks. Models such as Prithvi-100M \citep{jakubik2024prithvi}, SatMAE \citep{cong2022satmae}, and AlphaEarth Foundations \citep{brown2025alphaearth} learn rich representations from large-scale satellite imagery through self-supervised pre-training, enabling strong downstream performance with limited labeled data. However, a critical gap exists between how these models are pre-trained and how they are deployed.

Prithvi-100M, for example, is pre-trained exclusively on 6-band Harmonized Landsat Sentinel-2 (HLS) imagery from the contiguous United States. In practice, users need to apply such models to heterogeneous inputs: 4-band commercial optical imagery (e.g., GaoFen-2), 2-band SAR data (Sentinel-1 VV/VH), 10-band Sentinel-2, or arbitrary combinations thereof. The standard workaround is zero-padding---filling missing channels with zeros and truncating excess channels---but this approach is uninformed: it wastes channel capacity and provides no mechanism for the model to learn cross-modal mappings.

Parameter-efficient fine-tuning (PEFT) methods offer a promising path to adapt frozen foundation models to new tasks and domains without full retraining. Methods such as LoRA \citep{hu2022lora}, BitFit \citep{zaken2022bitfit}, and Houlsby Adapters \citep{houlsby2019adapter} have demonstrated strong results in NLP and computer vision. IBM's TerraTorch toolkit \citep{ibm2025terratorch} provides LoRA and Adapter support for Prithvi, and Esri's ArcGIS platform offers pre-packaged Prithvi models for specific tasks. Yet no systematic study has evaluated how these PEFT methods perform when the input modality differs from pre-training---the very scenario that practitioners face most often.

This paper addresses this gap through a comprehensive empirical study. We benchmark five PEFT methods on a frozen Prithvi-100M backbone across five modality configurations, ranging from the full 10-band Sentinel-2 input to extreme out-of-distribution 2-band inputs. Our 75-experiment evaluation (5 methods $\times$ 5 modalities $\times$ 3 seeds) on EuroSAT \citep{helber2019eurosat} reveals several findings that challenge common assumptions:

\begin{enumerate}
\item \textbf{LoRA fails on Prithvi.} Despite using 155K trainable parameters, LoRA produces results indistinguishable from a 7.7K-parameter Linear Probe ($\Delta$OA~$<$~0.001). We attribute this to Prithvi's fused QKV attention structure, which differs from the separate Q/K/V matrices that LoRA was designed for.

\item \textbf{Houlsby adapters dominate and enable cross-modal transfer.} With 1.2M parameters inserted as bottleneck layers after each Transformer block, Houlsby achieves +16--23\% OA improvement. Notably, it achieves higher accuracy on 4-channel rgb\_sar input (0.738) than on 3-channel rgb input (0.727), demonstrating that backbone-internal adapters can implicitly learn to utilize zero-padded channels.

\item \textbf{Input-stage adaptation does not work.} We propose GeoAdapter, a residual input-stage adapter with channel projection, squeeze-excitation attention, and spatial refinement ($\sim$150 parameters). Despite its residual design guaranteeing a zero-pad lower bound at initialization, GeoAdapter consistently underperforms the baseline by 3--22\%, with unstable optimization dynamics.

\item \textbf{Modality selection dominates PEFT method selection.} The accuracy gap between the best and worst modality configurations within a single method (up to 27\%) far exceeds the gap between the best and worst methods within a single modality (up to 16\%).
\end{enumerate}

We release our evaluation platform, Geo-MLOps, which provides a unified training engine, benchmark runner with YAML-driven experiment configurations, and a web-based monitoring dashboard. All code and experiment results are publicly available.

% ============================================================================
\section{Related Work}
\label{sec:related}
% ============================================================================

\subsection{Geospatial foundation models}

The emergence of large-scale self-supervised pre-training has produced several geospatial foundation models. Prithvi-100M \citep{jakubik2024prithvi} uses a temporal Vision Transformer (ViT) with masked autoencoder (MAE) pre-training on HLS data, producing 768-dimensional features from 6-band inputs. SatMAE \citep{cong2022satmae} extends MAE with temporal positional encoding for multi-spectral imagery. AlphaEarth Foundations \citep{brown2025alphaearth} takes a different approach with a multi-resolution Space-Time-Precision architecture trained on 9+ data sources, producing 64-dimensional embeddings constrained to a unit hypersphere. While AlphaEarth provides pre-computed global embeddings, its model weights are not publicly available, making Prithvi the primary open-source option for fine-tuning research.

\subsection{Parameter-efficient fine-tuning}

PEFT methods adapt pre-trained models by updating only a small subset of parameters. \textbf{Linear Probing} trains only a classification head while freezing the entire backbone. \textbf{BitFit} \citep{zaken2022bitfit} unfreezes all bias parameters, providing a lightweight adaptation signal throughout the network. \textbf{LoRA} \citep{hu2022lora} decomposes weight updates into low-rank matrices $\Delta W = AB$ where $A \in \mathbb{R}^{d \times r}$ and $B \in \mathbb{R}^{r \times d}$ with rank $r \ll d$, typically applied to attention projection matrices. \textbf{Houlsby Adapters} \citep{houlsby2019adapter} insert bottleneck modules (down-projection, nonlinearity, up-projection) after each Transformer layer, with zero-initialized up-projection ensuring identity initialization. These methods have been extensively validated in NLP and natural image classification, but their behavior on geospatial ViTs with heterogeneous spectral inputs remains unexplored.

\subsection{Geospatial ML frameworks}

Several frameworks address different layers of the geospatial ML stack. TorchGeo \citep{stewart2024torchgeo} provides standardized dataset loaders and samplers for geospatial data. eo-learn (Sentinel Hub) offers modular EOPatch/EOTask/EOWorkflow abstractions for satellite data processing pipelines. TerraTorch \citep{ibm2025terratorch} provides LoRA and Adapter support specifically for Prithvi fine-tuning. However, none of these frameworks include systematic benchmarks evaluating PEFT methods under modality mismatch conditions---the scenario we address in this work.

% ============================================================================
\section{Method}
\label{sec:method}
% ============================================================================

\subsection{Backbone: Prithvi-100M}

We use Prithvi-100M \citep{jakubik2024prithvi} as the frozen backbone. The model consists of a Conv3d patch embedding layer (kernel size $1 \times 16 \times 16$, projecting 6 input channels to 768 dimensions), 12 Transformer encoder layers (hidden dimension 768, 12 attention heads, MLP ratio 4, GELU activation), and a final LayerNorm. For single-timestep inputs, we squeeze the temporal dimension of the Conv3d weight tensor, yielding an equivalent Conv2d($6, 768, 16, 16$). We load all 149 weight tensors from the official checkpoint, mapping Prithvi's key format (\texttt{encoder.blocks.N.attn.qkv}) to PyTorch's standard TransformerEncoderLayer format (\texttt{blocks.N.self\_attn.in\_proj\_weight}). All backbone parameters are frozen; only PEFT-injected parameters and the task head are trainable.

\subsection{PEFT methods}

We evaluate five methods spanning different adaptation strategies:

\textbf{Linear Probe.} Only a linear classification head ($768 \to C$ classes) is trained. The backbone and a zero-padding input adapter are completely frozen. This serves as the lower-bound baseline (7,690 trainable parameters for $C=10$).

\textbf{BitFit} \citep{zaken2022bitfit}. All bias parameters throughout the backbone are unfrozen while weights remain frozen, plus the classification head. This provides a distributed adaptation signal at minimal parameter cost (110,602 trainable parameters).

\textbf{LoRA} \citep{hu2022lora}. We inject low-rank adaptation matrices ($r=8$) into the self-attention layers of each Transformer block. For each target linear layer with weight $W_0 \in \mathbb{R}^{d \times d}$, we add $\Delta W = AB$ where $A \in \mathbb{R}^{d \times 8}$ and $B \in \mathbb{R}^{8 \times d}$, with $A$ initialized from $\mathcal{N}(0, 1/\sqrt{r})$ and $B$ initialized to zero (155,146 trainable parameters).

\textbf{Houlsby Adapter} \citep{houlsby2019adapter}. After each Transformer layer, we insert a bottleneck module: $\text{Adapter}(x) = x + W_{\text{up}} \cdot \text{GELU}(W_{\text{down}} \cdot x)$, where $W_{\text{down}} \in \mathbb{R}^{768 \times 64}$ and $W_{\text{up}} \in \mathbb{R}^{64 \times 768}$, with $W_{\text{up}}$ zero-initialized to ensure identity at initialization (1,197,322 trainable parameters across 12 layers).

\textbf{GeoAdapter (input-stage residual adapter).} We propose a lightweight adapter inserted \emph{before} the frozen patch embedding, designed to learn a residual correction on top of zero-padding:
\begin{equation}
\text{GeoAdapter}(x) = \text{ZeroPad}(x) + \alpha \cdot f(x)
\end{equation}
where $\alpha$ is a learnable scalar initialized to 0 (ensuring the initial output equals zero-padding), and $f$ consists of three layers: (1)~a $1 \times 1$ convolution projecting $C_{\text{in}}$ channels to 6; (2)~a squeeze-excitation channel attention module \citep{hu2022se} with reduction ratio 2; and (3)~a depthwise $3 \times 3$ convolution with batch normalization and GELU activation. Total adapter parameters: $\sim$150 (7,826 including the classification head).

\subsection{Modality configurations}

To simulate real-world heterogeneous input scenarios, we define five modality configurations by selecting different subsets of EuroSAT's 13 Sentinel-2 bands:

\begin{table}[htbp]
\centering
\caption{Modality configurations. All inputs are zero-padded (or truncated) to 6 channels to match Prithvi's expected input, except GeoAdapter which uses learned residual adaptation.}
\label{tab:modalities}
\begin{tabular}{llcl}
\toprule
ID & Bands selected & $C_{\text{in}}$ & Simulates \\
\midrule
s2\_full & B2--B12 (first 10) & 10 & Full Sentinel-2 (superset) \\
rgb & B4, B3, B2 & 3 & Visible-only sensor \\
rgb\_sar & B4, B3, B2, B1 & 4 & Optical + SAR mix \\
gf2 & B4, B3, B2, B8 & 4 & GaoFen-2 (B,G,R,NIR) \\
sar\_only & B1, B2 & 2 & SAR-only (extreme OOD) \\
\bottomrule
\end{tabular}
\end{table}

\subsection{Geo-MLOps platform}

All experiments are conducted through Geo-MLOps, an open-source platform we developed for reproducible PEFT evaluation. The platform consists of: (1)~a \texttt{geoadapter} Python package containing all adapter implementations, the Prithvi backbone loader, a unified training engine with differential learning rates and cosine annealing, and a YAML-driven benchmark runner; (2)~a FastAPI backend with WebSocket-based real-time training monitoring; and (3)~a Vue~3 dashboard for interactive experimentation. The \texttt{geoadapter} package has zero web framework dependencies and can be used standalone in Colab or any Python environment.

% ============================================================================
\section{Experiments}
\label{sec:experiments}
% ============================================================================

\subsection{Setup}

\textbf{Dataset.} EuroSAT \citep{helber2019eurosat}: 27,000 Sentinel-2 images at 64$\times$64 pixels, 13 spectral bands, 10 land-use classes. We use the standard train/test split provided by TorchGeo \citep{stewart2024torchgeo}.

\textbf{Training.} 50 epochs, batch size 64, AdamW optimizer (weight decay 0.01), cosine annealing learning rate schedule. Learning rate: $10^{-3}$ for classification head and input adapter; $10^{-4}$ for backbone-internal PEFT parameters (LoRA, BitFit, Houlsby). Per-channel z-score normalization applied to all inputs.

\textbf{Evaluation.} Overall Accuracy (OA) and Macro F1 on the test set. Each experiment is repeated with 3 random seeds (42, 123, 456); we report mean $\pm$ standard deviation.

\textbf{Hardware.} NVIDIA A100 40GB (Google Colab Pro+). Total compute: $\sim$27 GPU-hours for 75 experiments.

\subsection{Main results}

\begin{table}[htbp]
\centering
\caption{Overall Accuracy (mean $\pm$ std across 3 seeds) for 5 PEFT methods $\times$ 5 modality configurations on EuroSAT. Best result per column in \textbf{bold}.}
\label{tab:main}
\begin{tabular}{lcccccc}
\toprule
Method & s2\_full & rgb & rgb\_sar & gf2 & sar\_only & Params \\
\midrule
Linear Probe & .657{\scriptsize$\pm$.002} & .556{\scriptsize$\pm$.003} & .552{\scriptsize$\pm$.002} & .649{\scriptsize$\pm$.005} & .387{\scriptsize$\pm$.003} & 7.7K \\
BitFit & .702{\scriptsize$\pm$.002} & .608{\scriptsize$\pm$.001} & .605{\scriptsize$\pm$.001} & .701{\scriptsize$\pm$.001} & .449{\scriptsize$\pm$.002} & 111K \\
LoRA ($r$=8) & .658{\scriptsize$\pm$.002} & .556{\scriptsize$\pm$.002} & .553{\scriptsize$\pm$.002} & .650{\scriptsize$\pm$.005} & .388{\scriptsize$\pm$.004} & 155K \\
Houlsby ($d$=64) & \textbf{.821}{\scriptsize$\pm$.002} & \textbf{.727}{\scriptsize$\pm$.002} & \textbf{.738}{\scriptsize$\pm$.001} & \textbf{.820}{\scriptsize$\pm$.002} & \textbf{.615}{\scriptsize$\pm$.003} & 1.2M \\
GeoAdapter & .535{\scriptsize$\pm$.031} & .332{\scriptsize$\pm$.021} & .432{\scriptsize$\pm$.058} & .445{\scriptsize$\pm$.070} & .356{\scriptsize$\pm$.007} & 7.8K \\
\bottomrule
\end{tabular}
\end{table}

\begin{table}[htbp]
\centering
\caption{Improvement in OA over Linear Probe baseline ($\Delta$OA). Positive values indicate improvement.}
\label{tab:delta}
\begin{tabular}{lccccc}
\toprule
Method & s2\_full & rgb & rgb\_sar & gf2 & sar\_only \\
\midrule
BitFit & +.045 & +.052 & +.053 & +.052 & +.062 \\
LoRA ($r$=8) & +.001 & +.000 & +.001 & +.001 & +.001 \\
Houlsby ($d$=64) & \textbf{+.164} & \textbf{+.171} & \textbf{+.185} & \textbf{+.171} & \textbf{+.228} \\
GeoAdapter & $-$.122 & $-$.224 & $-$.120 & $-$.204 & $-$.031 \\
\bottomrule
\end{tabular}
\end{table}

\subsection{Finding 1: LoRA fails on Prithvi-100M}

LoRA with rank 8 and 155K trainable parameters produces results statistically indistinguishable from Linear Probe (7.7K parameters) across all five modality configurations. The maximum $\Delta$OA is +0.001 (Table~\ref{tab:delta}).

We attribute this to Prithvi's attention architecture. Prithvi uses a fused QKV projection (\texttt{attn.qkv.weight} $\in \mathbb{R}^{2304 \times 768}$) rather than separate Q, K, V matrices. When mapped to PyTorch's standard \texttt{self\_attn.in\_proj\_weight}, LoRA's low-rank decomposition operates on the concatenated QKV space. This may prevent LoRA from selectively modifying query-key interactions, which is the mechanism through which LoRA typically achieves adaptation in standard ViTs.

This finding has practical implications: practitioners using TerraTorch or similar toolkits should not assume that LoRA---the default PEFT method in many frameworks---will be effective on Prithvi without verifying the attention structure compatibility.

\subsection{Finding 2: Houlsby adapters dominate and enable cross-modal transfer}

Houlsby adapters achieve the highest accuracy across all modality configurations, with improvements of +16.4\% to +22.8\% over Linear Probe (Table~\ref{tab:delta}). The improvement is largest on sar\_only (+22.8\%), the most out-of-distribution configuration.

A particularly noteworthy result is that Houlsby on rgb\_sar (0.738) \emph{outperforms} Houlsby on rgb (0.727) by +1.1\%. Since rgb\_sar differs from rgb only by an additional zero-padded channel, this demonstrates that Houlsby's bottleneck adapters can learn to extract useful signal from channels that contain no explicit information at input time. In contrast, for Linear Probe and BitFit, rgb\_sar and rgb produce nearly identical results ($\Delta < 0.003$), confirming that without sufficient backbone adaptation capacity, zero-padded channels are effectively ignored.

This suggests that the ``modality gap'' problem may not require specialized input-stage adaptation. Instead, backbone-internal adapters with sufficient capacity can learn cross-modal representations through the Transformer's attention mechanism, even when the additional modality is presented as zero-padded channels.

\subsection{Finding 3: Input-stage adaptation does not work}

GeoAdapter, our proposed residual input-stage adapter, consistently underperforms the Linear Probe baseline by 3--22\% across all modality configurations (Table~\ref{tab:delta}). This is a negative result that we report in full for the benefit of the community.

The residual scale parameter $\alpha$, initialized to 0 to ensure the initial output equals zero-padding, rapidly grows to $|\alpha| > 1.0$ during training. Across modalities, the mean absolute scale ranges from 1.18 (s2\_full) to 1.82 (sar\_only), with inconsistent sign across seeds (e.g., rgb\_sar: $\alpha \in \{-1.40, +1.51, -1.37\}$). This indicates severe optimization instability: the adapter's learned signal overwhelms the zero-pad baseline rather than refining it.

The root cause is a parameter-gradient mismatch. The adapter has only $\sim$150 learnable parameters, but receives gradient signal from the 7,690-parameter classification head backpropagated through the 86M-parameter frozen backbone. This creates a noisy, high-variance optimization landscape where the tiny adapter cannot converge to a stable solution.

\subsection{Finding 4: Modality selection dominates PEFT method selection}

The accuracy gap between the best and worst modality configurations within a single PEFT method is consistently larger than the gap between the best and worst methods within a single modality:

\begin{itemize}
\item Within Linear Probe: s2\_full (0.657) vs.\ sar\_only (0.387) = 0.270 gap
\item Within Houlsby: s2\_full (0.821) vs.\ sar\_only (0.615) = 0.206 gap
\item Within s2\_full: Houlsby (0.821) vs.\ Linear Probe (0.657) = 0.164 gap
\end{itemize}

This has a direct practical implication: for practitioners with limited compute budgets, investing in better input data (more spectral bands, sensors closer to the pre-training distribution) yields larger returns than investing in more sophisticated PEFT methods.

% ============================================================================
\section{Discussion}
\label{sec:discussion}
% ============================================================================

\textbf{Why does LoRA fail while Houlsby succeeds?} Both methods modify the Transformer backbone, but at different granularities. LoRA modifies the attention computation through low-rank weight perturbations, while Houlsby inserts entirely new computational pathways (bottleneck modules) after each layer. Our results suggest that for Prithvi's MAE-pretrained features, the attention patterns are already well-suited for feature extraction---the bottleneck lies in the \emph{feature transformation} between layers, which Houlsby directly addresses. This aligns with recent findings that adapter-based methods outperform LoRA on vision tasks when the domain shift is large.

\textbf{Why does input-stage adaptation fail?} Our GeoAdapter experiment provides a clear negative result: a lightweight adapter ($\sim$150 parameters) inserted before the frozen backbone cannot learn a useful modality mapping in 50 epochs. The fundamental issue is that the adapter must simultaneously learn two things: (1)~how to map arbitrary input channels to the pre-training distribution, and (2)~how to do so in a way that is useful for the downstream task. With the backbone frozen, the only gradient signal comes from the classification head, which is too indirect to guide the adapter's channel mapping. A potential remedy would be to combine input-stage adaptation with backbone-internal PEFT (e.g., GeoAdapter + Houlsby), but this would conflate the contributions and is left for future work.

\textbf{Implications for the Houlsby cross-modal finding.} The observation that Houlsby on rgb\_sar outperforms Houlsby on rgb suggests that zero-padded channels are not ``wasted'' when the backbone has sufficient adaptation capacity. The Transformer's self-attention mechanism can learn to attend to spatial patterns in the zero-padded channels that correlate with useful features in the non-zero channels. This has implications for multi-sensor fusion: rather than designing complex input-stage fusion modules, practitioners may achieve better results by simply concatenating and zero-padding heterogeneous inputs and relying on backbone-internal adapters to learn the fusion implicitly.

\textbf{Limitations.} Our study has several limitations. (1)~We evaluate on a single dataset (EuroSAT) with a single task (scene classification). Extending to segmentation and multi-label tasks would strengthen the findings. (2)~Our LoRA implementation operates on PyTorch's standard TransformerEncoderLayer rather than Prithvi's native architecture; a native implementation might yield different results. (3)~The modality configurations use band subsets from the same sensor (Sentinel-2) rather than truly heterogeneous multi-sensor data. (4)~We use a single backbone (Prithvi-100M); results may differ for other GeoFMs such as SatMAE or Clay.

% ============================================================================
\section{Conclusion}
\label{sec:conclusion}
% ============================================================================

We presented the first systematic benchmark of PEFT methods for adapting geospatial foundation models to heterogeneous spectral inputs. Through 75 experiments on Prithvi-100M across 5 PEFT strategies and 5 modality configurations, we established four findings with practical implications:

\begin{enumerate}
\item LoRA, the most popular PEFT method, fails on Prithvi's fused-QKV architecture---practitioners should verify attention structure compatibility before adopting LoRA for geospatial models.
\item Houlsby adapters are the most effective strategy, achieving +16--23\% OA improvement and demonstrating implicit cross-modal transfer through zero-padded channels.
\item Input-stage modality adaptation with a lightweight residual adapter does not work---the modality gap is better addressed through backbone-internal adaptation.
\item Modality selection (which bands to use) has a larger impact on accuracy than PEFT method selection (how to fine-tune).
\end{enumerate}

These findings provide actionable guidance for deploying geospatial foundation models in real-world scenarios where input data does not match pre-training specifications. Our open-source Geo-MLOps platform enables reproducible evaluation of PEFT methods on any geospatial foundation model.

% ============================================================================
% References
% ============================================================================

\begin{thebibliography}{99}

\bibitem[Jakubik et~al.(2024)]{jakubik2024prithvi}
Jakubik, J., Roy, S., Phillips, C.~E., et~al.
\newblock Foundation models for generalist geospatial artificial intelligence.
\newblock {\em Nature}, 2024.

\bibitem[Brown et~al.(2025)]{brown2025alphaearth}
Brown, C.~F., Kazmierski, M.~R., Pasquarella, V.~J., et~al.
\newblock {AlphaEarth Foundations}: An embedding field model for accurate and efficient global mapping from sparse label data.
\newblock {\em arXiv preprint arXiv:2507.22291}, 2025.

\bibitem[Hu et~al.(2022)]{hu2022lora}
Hu, E.~J., Shen, Y., Wallis, P., et~al.
\newblock {LoRA}: Low-rank adaptation of large language models.
\newblock In {\em ICLR}, 2022.

\bibitem[Houlsby et~al.(2019)]{houlsby2019adapter}
Houlsby, N., Giurgiu, A., Jastrzebski, S., et~al.
\newblock Parameter-efficient transfer learning for {NLP}.
\newblock In {\em ICML}, 2019.

\bibitem[Zaken et~al.(2022)]{zaken2022bitfit}
Zaken, E.~B., Goldberg, Y., and Ravfogel, S.
\newblock {BitFit}: Simple parameter-efficient fine-tuning for transformer-based masked language-models.
\newblock In {\em ACL}, 2022.

\bibitem[Cong et~al.(2022)]{cong2022satmae}
Cong, Y., Khanna, S., Meng, C., et~al.
\newblock {SatMAE}: Pre-training transformers for temporal and multi-spectral satellite imagery.
\newblock In {\em NeurIPS}, 2022.

\bibitem[Helber et~al.(2019)]{helber2019eurosat}
Helber, P., Bischke, B., Dengel, A., and Borth, D.
\newblock {EuroSAT}: A novel dataset and deep learning benchmark for land use and land cover classification.
\newblock {\em IEEE J-STARS}, 12(7):2217--2226, 2019.

\bibitem[Stewart et~al.(2024)]{stewart2024torchgeo}
Stewart, A.~J., Robinson, C., Corley, I.~A., et~al.
\newblock {TorchGeo}: Deep learning with geospatial data.
\newblock In {\em ACM SIGSPATIAL}, 2024.

\bibitem[IBM(2025)]{ibm2025terratorch}
IBM Research.
\newblock Simplifying geospatial {AI} with {TerraTorch} 1.0.
\newblock \url{https://research.ibm.com/blog/simplifying-geospatial-ai-with-terra-torch-1-0}, 2025.

\bibitem[Hu and Rosser(2022)]{hu2022se}
Hu, J., Shen, L., and Sun, G.
\newblock Squeeze-and-excitation networks.
\newblock {\em IEEE TPAMI}, 42(8):2011--2023, 2020.

\end{thebibliography}

\end{document}
